% ---------------------------------------------------------------------------
% Gritt - Robotics Controls Intern
% Tailored from bases/simulation_robotics.tex on 2026-08-25
% Worksheet: notes.md in this folder. Posting: job_description.txt
%
% Do NOT put the company name anywhere in the rendered document. It belongs in
% the filename only.
% ---------------------------------------------------------------------------
\documentclass[10pt,letterpaper]{article}
\usepackage{resume}

% Tighter entry rhythm for a dense one-page technical resume.
\setlist[itemize]{topsep=1pt,itemsep=0.5pt,after=\vspace{1pt}}
\renewcommand{\entryspace}{\vspace{0pt}}

\begin{document}
\bodyfont

\resumeheader
\educationsection

\sectionheader{TECHNICAL SKILLS}

\techline{\textbf{Languages}: C/C++, Python, Java}
\techline{\textbf{Controls \& Dynamics}: Fixed-point PI/PID, saturation, anti-windup, slew limiting, state estimation, sensor fusion, Bode/Nyquist}
\techline{\textbf{Robotics \& Embedded}: PX4, visual-inertial odometry, IMU fusion, STM32H7, CAN, real-time firmware}
\techline{\textbf{Simulation \& Validation}: Software-in-the-loop, closed-loop simulation, scripted replay, regression tests, parameter sweeps, NumPy, pandas, Matplotlib}

\sectionheader{EXPERIENCE}

\roleline{Software Lead}{Jun 2026 - Present}
\orgline{\spartanracing}
\begin{itemize}
  \item Lead \textbf{12 engineers} across \textbf{10+ vehicle-control, modeling, validation, and telemetry projects} for a Formula SAE electric racecar
  \item Developing the custom-inverter VCU path for a \textbf{four-motor} EV, mapping calibrated pedal input to bounded CAN current commands on a \textbf{10 ms} loop with fail-safe precharge and ready-to-drive gating
  \item Lead battery, BMS, motor, and inverter plant-model integration for a core team of \textbf{7 engineers}, defining interfaces so VCU control code can be evaluated before hardware tests
\end{itemize}

\entryspace

\roleline{Software Engineer Intern}{Feb 2026 - Present}
\orgline{Argus Defense \textbar{} San Jose, CA}
\begin{itemize}
  \item Tuned PX4 PID flight-control loops to reduce vertical and horizontal drift during autonomous position hold
  \item Integrated and tuned visual-inertial odometry for onboard state estimation, connecting perception output to the flight-control stack and evaluating changes against logged telemetry
  \item Developed embedded flight-control changes in C and iterated gains through repeatable hardware tests and log review
\end{itemize}

\entryspace

\roleline{Software Controls Engineer}{Jul 2025 - Jun 2026}
\orgline{\spartanracing}
\begin{itemize}
  \item Developed fixed-point PI firmware with saturation, anti-windup, and slew limiting that held measured power within \textbf{0.01\%} of its commanded setpoint while enforcing an \textbf{80 kW} limit
  \item Built Python telemetry parsers and a software-in-the-loop replay harness for recorded CSV logs, controller parameter sweeps, and regression testing
  \item Planned and executed on-vehicle tests, measured control response and latency, and generated plots for tuning and baseline comparison
\end{itemize}

\entryspace

\roleline{Software Intern}{Aug 2024 - Jun 2025}
\orgline{\spartanracing}
\begin{itemize}
  \item Tuned the PI torque controller to within \textbf{2\%} error using logged step responses, and the vehicle placed \textbf{1st in endurance} at MIS EV 2025
\end{itemize}

\sectionheader{PROJECT EXPERIENCE}

\roleline{SR-Wireless-CAN, Control-System Test Platform}{2026}
\orgline{\spartanracing \textbar{} Backend lead}
\techline{Python, CAN, recorded-trace replay, FastAPI, WebSockets, SQLite, Docker}
\begin{itemize}
  \item Led backend development for a test platform supporting \textbf{3 capture modes}, including simulated data, recorded-trace replay, and live CAN, plus \textbf{3 signal-monitoring views}
  \item Reverse-engineered and implemented an \textbf{8-stage} CAN firmware-deployment workflow from vendor traces, with acknowledgement checks, CRC verification, and persisted run history
\end{itemize}

\end{document}

